\section{Discussion and extensions}

The separation comparisons provide two benchmarks for the sensitivity of homogeneity testing: one varies the outcome envelope at matched primitive smoothness, and the other varies the assignment information available to the test on a shared model class. The first comparison, in \cref{obj:thm:heavy-tail-comparison}, measures the cost of allowing outcomes subject to a finite conditional moment bound relative to bounded outcomes. The second, in \cref{obj:thm:exact-propensity-preservation-boundary}, measures the value of supplying the true propensity function. Together, they characterize how outcome tails and assignment information affect the smallest detectable departure from a constant conditional mean effect.

Matching primitive smoothness makes the outcome comparison interpretable. The propensity and baseline smoothness exponents describe nuisance complexity, whereas the effect smoothness exponent describes the complexity of the departure being tested. The bounded benchmark in \cref{obj:def:bounded-model} retains these three restrictions. Consequently, the radius ratio in \cref{obj:def:comparison-handle} compares outcome envelopes while keeping the stipulated effect and nuisance complexity fixed. The regions in \cref{obj:def:tail-region,obj:def:equality-region} distinguish a diverging separation penalty from a uniformly bounded radius ratio. These are benchmarks for worst-case sensitivity over the specified classes: they identify when the conditional moment envelope changes the order of the testing problem.

The assignment comparison separates the complexity of the effect from the information needed to handle the nuisance functions. On the shared original class, \cref{obj:thm:exact-propensity-preservation-boundary} identifies the exact region where the unknown-propensity experiment attains the supplied-propensity separation order. The broader-class result in \cref{obj:thm:oracle-frontier-broad-class} gives a complementary interpretation. Its model, defined in \cref{obj:def:oracle-broad-model}, permits continuous propensity and baseline mean functions subject to the retained overlap and mean envelopes, together with the conditional moment and effect smoothness restrictions. Attainment over this class supports a broader nuisance scope when the true propensity is supplied. The resulting separation order is governed by the moment order and effect smoothness, clarifying the role of assignment information in accommodating nuisance variation.

The population target also has a direct variance interpretation. Writing the distance from a constant effect as
\[
d(P)=
\left[
\int_0^1
\left\{\tau_P(x)-\int_0^1\tau_P(z)\,dz\right\}^{2}\,dx
\right]^{1/2},
\]
the uniform covariate design in \cref{obj:ass:uniform} gives
\[
d(P)^2=\operatorname{Var}_{P}\{\tau_P(X)\}.
\]
Thus separation is measured in the standard deviation of the conditional mean treatment effect across covariate values. Under the causal completion in \cref{obj:def:causal-completion}, this contrast equals the conditional expectation of the potential-outcome difference. The target therefore quantifies heterogeneity in conditional mean effects, with the covariate distribution determining the population weighting.

Public finite-sample quantities connect these comparisons to calibrated rejection rules. The construction in \cref{obj:def:explicit-score-ledger} specifies the projection ranks, clipping levels, bias and covariance bounds, and rejection cutoff. Whole-null calibration in \cref{obj:thm:whole-null-calibration-resolved} accounts for the unknown constant effect through profiling. Combined with the attainment guarantees in \cref{obj:thm:heavy-tail-comparison}, the public sample-size thresholds identify when the stated power separation lies inside the model's nonempty alternative domain. The analogous supplied-propensity guarantees appear in \cref{obj:thm:oracle-frontier-broad-class,obj:thm:oracle-frontier-resolved}.

The explicit sufficient constants are deliberately conservative. The following values come directly from the displayed formulas; logarithms make their scale readable. Here \(N_v^{\mathrm{att}}\) is the original-record nonempty-power threshold and \(N_v^{\mathrm e}\) is its supplied-propensity counterpart.
\begin{center}
\begin{tabular}{llllll}
\hline
\(p\) & \((\alpha,\beta,\gamma)\) & regime & \(E\) & \(\log_{10}C_v^{\mathrm{att}}\) & \(\log_{10}N_v^{\mathrm{att}}\; /\;\log_{10}N_v^{\mathrm e}\) \\
\hline
\(2\) & \((1,1,1)\) & effect & \(0.400\) & \(8.785\) & \(25.35\;/\;12.86\) \\
\(3/2\) & \((1,1,1)\) & effect & \(0.286\) & \(8.786\) & \(35.49\;/\;18.00\) \\
\(3/2\) & \((1/20,1/20,1)\) & rough & \(0.154\) & \(9.934\) & \(73.37\;/\;18.00\) \\
\hline
\end{tabular}
\end{center}
For the first row, for example, the supplied-propensity formula gives \(N_v^{\mathrm e}=7{,}183{,}754{,}725{,}838\). These are analytic sufficient thresholds for the stated finite-sample power guarantee, rather than estimates of the sample size needed in practice. The theorem's main quantitative content is the minimax power of \(n\) and the matching lower bound; practical calibration of the constants would require implementation and simulation beyond the present analysis.

\subsection{Limitations and future work}

The analysis treats a known uniform covariate design on the unit interval and public moment and smoothness exponents. Extending the construction to an unknown covariate distribution would require handling the population weighting and the estimation of projection geometry together. Higher-dimensional covariates raise a related question about how approximation complexity and nuisance smoothness interact in the separation comparison.

Adaptive selection of the moment and smoothness exponents is another research direction. A useful extension would quantify the cost of choosing clipping levels and projection ranks from the data while maintaining calibration over the entire constant-effect null and establishing nonempty power domains.

Finally, the supplied-function experiment provides a benchmark for studying estimated assignment functions. Such an extension would quantify how propensity approximation error and the dependence introduced by its estimation enter the bias bounds, rejection threshold, and attainable separation. Comparing those guarantees with the exact supplied-function benchmark would clarify the assignment accuracy needed to retain its sensitivity to effect heterogeneity.
